% GENERATED FILE. Edit canonical lessons, then run npm run generate:books.
% canonical-source-hash: fnv1a64:7eec08aabb49b358
% canonical-lessons: PA-C150-siuna, PA-C150-bharna, PA-C150-khiccna, PA-C150-jagna, PA-C150-calauna

\chapter{To sew, To fill, To pull, To wake up, To drive}
\label{ch:pa-to-sew-to-fill-to-pull-to-wake-up-to-drive}
\hlchaptermodality{150}

\begin{chapteropening}
\textbf{By the end of this chapter:} \emph{I can say to sew, to fill, to pull, to wake up and to drive.}

Complete the last lesson of chapter 150: I can say to sew, to fill, to pull, to wake up and to drive.
\end{chapteropening}

\section[\texorpdfstring{\pa{ਸਿਉਣਾ} (siuṇā)}{ਸਿਉਣਾ (siuṇā)}]{\pa{ਸਿਉਣਾ} (siuṇā) — to sew}

\label{lesson:PA-C150-siuna}



\begin{quote}
Before the new one: say the Punjabi for to use, then the Punjabi for to make.
\end{quote}

\subsection*{You'll want to know: \pa{ਸਿਉਣਾ}}
\textbf{\pa{ਸਿਉਣਾ}} — \emph{siuṇā} — \textquotedblleft{}to sew\textquotedblright{}.

\begin{grammarlens}[title={What you've built}]
One more word for this chapter.
\end{grammarlens}

\subsection*{Your turn}
\begin{itemize}
\raggedright
  \item \emph{Say it:} \emph{siuṇā}
  \item \emph{Say it:} \emph{siuṇā}, once more
  \item \emph{Say it:} say \emph{siuṇā}
  \item \emph{Recall:} say the Punjabi for to use, then the Punjabi for to make, then say \emph{siuṇā} again
\end{itemize}

\begin{tcolorbox}[breakable,colback=teal!4,colframe=teal!35!black,title={Before you move on}]
What does \pa{ਸਿਉਣਾ} mean? (\textquotedblleft{}To sew\textquotedblright{}.) Say it once more.
\end{tcolorbox}
\section[\texorpdfstring{\pa{ਭਰਨਾ} (bharnā)}{ਭਰਨਾ (bharnā)}]{\pa{ਭਰਨਾ} (bharnā) — to fill}

\label{lesson:PA-C150-bharna}



\begin{quote}
Before the new one: say the Punjabi for to make, then the Punjabi for to sew.
\end{quote}

\subsection*{You'll want to know: \pa{ਭਰਨਾ}}
\textbf{\pa{ਭਰਨਾ}} — \emph{bharnā} — \textquotedblleft{}to fill\textquotedblright{}.

\begin{grammarlens}[title={What you've built}]
One more word for this chapter.
\end{grammarlens}

\subsection*{Your turn}
\begin{itemize}
\raggedright
  \item \emph{Say it:} \emph{bharnā}
  \item \emph{Say it:} \emph{bharnā}, once more
  \item \emph{Say it:} say \emph{bharnā}
  \item \emph{Recall:} say the Punjabi for to make, then the Punjabi for to sew, then say \emph{bharnā} again
\end{itemize}

\begin{tcolorbox}[breakable,colback=teal!4,colframe=teal!35!black,title={Before you move on}]
What does \pa{ਭਰਨਾ} mean? (\textquotedblleft{}To fill\textquotedblright{}.) Say it once more.
\end{tcolorbox}
\section[\texorpdfstring{\pa{ਖਿੱਚਣਾ} (khiccṇā)}{ਖਿੱਚਣਾ (khiccṇā)}]{\pa{ਖਿੱਚਣਾ} (khiccṇā) — to pull}

\label{lesson:PA-C150-khiccna}



\begin{quote}
Before the new one: say the Punjabi for to sew, then the Punjabi for to fill.
\end{quote}

\subsection*{You'll want to know: \pa{ਖਿੱਚਣਾ}}
\textbf{\pa{ਖਿੱਚਣਾ}} — \emph{khiccṇā} — \textquotedblleft{}to pull\textquotedblright{}.

\begin{grammarlens}[title={What you've built}]
One more word for this chapter.
\end{grammarlens}

\subsection*{Your turn}
\begin{itemize}
\raggedright
  \item \emph{Say it:} \emph{khiccṇā}
  \item \emph{Say it:} \emph{khiccṇā}, once more
  \item \emph{Say it:} say \emph{khiccṇā}
  \item \emph{Recall:} say the Punjabi for to sew, then the Punjabi for to fill, then say \emph{khiccṇā} again
\end{itemize}

\begin{tcolorbox}[breakable,colback=teal!4,colframe=teal!35!black,title={Before you move on}]
What does \pa{ਖਿੱਚਣਾ} mean? (\textquotedblleft{}To pull\textquotedblright{}.) Say it once more.
\end{tcolorbox}
\section[\texorpdfstring{\pa{ਜਾਗਣਾ} (jāgṇā)}{ਜਾਗਣਾ (jāgṇā)}]{\pa{ਜਾਗਣਾ} (jāgṇā) — to wake up}

\label{lesson:PA-C150-jagna}



\begin{quote}
Before the new one: say the Punjabi for to fill, then the Punjabi for to pull.
\end{quote}

\subsection*{You'll want to know: \pa{ਜਾਗਣਾ}}
\textbf{\pa{ਜਾਗਣਾ}} — \emph{jāgṇā} — \textquotedblleft{}to wake up\textquotedblright{}.

\begin{grammarlens}[title={What you've built}]
One more word for this chapter.
\end{grammarlens}

\subsection*{Your turn}
\begin{itemize}
\raggedright
  \item \emph{Say it:} \emph{jāgṇā}
  \item \emph{Say it:} \emph{jāgṇā}, once more
  \item \emph{Say it:} say \emph{jāgṇā}
  \item \emph{Recall:} say the Punjabi for to fill, then the Punjabi for to pull, then say \emph{jāgṇā} again
\end{itemize}

\begin{tcolorbox}[breakable,colback=teal!4,colframe=teal!35!black,title={Before you move on}]
What does \pa{ਜਾਗਣਾ} mean? (\textquotedblleft{}To wake up\textquotedblright{}.) Say it once more.
\end{tcolorbox}
\section[\texorpdfstring{\pa{ਚਲਾਉਣਾ} (calāuṇā)}{ਚਲਾਉਣਾ (calāuṇā)}]{\pa{ਚਲਾਉਣਾ} (calāuṇā) — to drive, to run (a machine)}

\label{lesson:PA-C150-calauna}



\begin{quote}
Before the new one: say the Punjabi for to pull, then the Punjabi for to wake up.
\end{quote}

\subsection*{You'll want to know: \pa{ਚਲਾਉਣਾ}}
\textbf{\pa{ਚਲਾਉਣਾ}} — \emph{calāuṇā} — \textquotedblleft{}to drive, to run (a machine)\textquotedblright{}.

\begin{grammarlens}[title={What you've built}]
One more word for this chapter.
\end{grammarlens}

\subsection*{Your turn}
\begin{itemize}
\raggedright
  \item \emph{Say it:} \emph{calāuṇā}
  \item \emph{Say it:} \emph{calāuṇā}, once more
  \item \emph{Say it:} say \emph{calāuṇā}
  \item \emph{Recall:} say the Punjabi for to pull, then the Punjabi for to wake up, then say \emph{calāuṇā} again
\end{itemize}

\begin{tcolorbox}[breakable,colback=teal!4,colframe=teal!35!black,title={Before you move on}]
What does \pa{ਚਲਾਉਣਾ} mean? (\textquotedblleft{}To drive, to run (a machine)\textquotedblright{}.) Say it once more.
\end{tcolorbox}
